\section{正则化、数值实现与估计不确定性}
\label{ch:lower-multifactor-estimation}

\begin{itemize}
  % Generated from curriculum/manifest.json; do not edit by hand.
\item 直接先修：下册《多因子模型：预测回归与风险价格》。

\end{itemize}

两个几乎相同的因子可能给出稳定的预测，却给出剧烈变化的系数。正则化针对的正是这种可识别方向不足的问题；先看目标怎样变化，再讨论参数选择。


\MFQLead{从正规方程到数值求解}


以下回归使用决策时点的信息集 $\mathcal I_t$。标准化等变换由当时已知的数据估计，因此模型的预测对应一个明确的时间起点。

线性模型 $y=X\beta+\varepsilon$ 的 OLS 目标是
\[
\min_\beta \frac12\lVert y-X\beta\rVert_2^2.
\]
满列秩时正规方程给出 $(X^{\mathsf T}X)\hat\beta=X^{\mathsf T}y$，但实现不应显式计算逆矩阵。QR 保留最小二乘几何，SVD 还能揭示近零奇异值；条件数检查回答“该样本能否稳定区分这些暴露”，不能被“求解器返回成功”替代。

面板研究还要声明误差依赖：同一资产跨期相关、同一期资产共同冲击、行业聚类与重叠持有期会改变标准误。固定效应吸收常数异质性，但不会自动解决未来信息、动态偏误或错误聚类层级。

\MFQLead{Ridge：用偏差换稳定}

Ridge 目标为
\[
\min_\beta \frac12\lVert y-X\beta\rVert_2^2+\frac\lambda2\lVert\beta\rVert_2^2,
\qquad
\hat\beta_{\mathrm{ridge}}=(X^{\mathsf T}X+\lambda I)^{-1}X^{\mathsf T}y.
\]
$\lambda>0$ 抬高小特征值，缓和共线性导致的方差放大，但引入偏差\cite{hoerl1970ridge}。惩罚参数必须只在训练/选择窗口决定；用最终测试期选 $\lambda$ 会把样本外变成调参样本。

设 $X=UDV^T$，奇异值为 $d_j$。把 $\beta=V\theta$ 代入 ridge 正规方程，逐坐标得到
\[
 \hat\theta_j=\frac{d_j}{d_j^2+\lambda}u_j^Ty.
\]
OLS 在该方向的系数为 $u_j^Ty/d_j$，所以 ridge 把它乘以 $d_j^2/(d_j^2+\lambda)$。小奇异值对应的方向收缩最多；它不是把所有系数乘同一常数。$d_1=10,d_2=0.1,\lambda=1$ 时两个收缩因子分别约为 $0.9901,0.0099$。

\MFQLead{Lasso：稀疏不是稳定选择}

Lasso 目标为
\[
\min_\beta \frac12\lVert y-X\beta\rVert_2^2+\lambda\lVert\beta\rVert_1.
\]
坐标下降更新第 $j$ 维时，对去除本维后的残差计算得分 $z_j=x_j^{\mathsf T}r_{-j}$，再使用软阈值
\[
\beta_j\leftarrow\frac{\operatorname{sign}(z_j)(|z_j|-\lambda)_+}{x_j^{\mathsf T}x_j}.
\]
该软阈值来自绝对值在零点的次梯度区间\cite{tibshirani1996regression}。零范数列不影响拟合，惩罚参数为正时可直接令其系数为零；高度相关特征中“选中哪一个”可能随小扰动改变，所以稀疏不等于经济解释稳定。

\MFQLead{同一组数据比较三种估计}

一列设计 $x=(1,2)$、$y=(1,1)$，有 $x^Tx=5,x^Ty=3$。OLS 为 $3/5$；ridge 取 $\lambda=1$ 得 $3/6=1/2$；lasso 的次梯度方程给 $\operatorname{soft}(3,1)/5=2/5$。若损失改为平均平方和 $\|y-X\beta\|^2/(2n)$，保持同一解需要把惩罚参数也除以 $n$。notebook 会显示这一列算例和两个共线特征的惩罚路径，允许改变噪声和惩罚强度。

\MFQTransition{滚动、扩展与加权估计}

扩展窗口使用全部历史，方差较小但对制度变化反应慢；固定滚动窗口遗忘旧数据，方差更大。指数权重最小化
\[
\sum_{s\le t}\lambda^{t-s}(y_s-x_s^{\mathsf T}\beta)^2,
\qquad 0<\lambda<1,
\]
在两者之间连续折中。半衰期 $h_{1/2}=\log(1/2)/\log\lambda$ 应用时间单位解释，而不是只报 $\lambda$。

横截面回归常按市值或误差方差加权。权重改变估计目标：市值权重更接近大盘可投资组合，等权更强调小股票，逆方差权重依赖一个额外噪声模型。权重若由未来成交量或事后波动估计，同样泄漏。

\MFQTransition{共线性、测量误差与稳定选择}

高度相关因子使 $X^{\mathsf T}X$ 条件数升高，单个系数对样本扰动敏感。Ridge 在相关方向上收缩系数，Lasso 则可能只保留其中一个变量。对样本作重抽样或改变估计窗口，可以观察选择频率与系数分布，了解这种选择有多稳定。

若特征带测量误差 $x^*=x+u$，简单回归系数通常向零衰减。财务数据库的回填、供应商映射错误和估计型特征都可能造成测量误差。正则化并没有建立测量误差模型，因此不能由此消除衰减偏差。

\MFQTransition{面板标准误与聚类层级}

误差可能在同一日期跨资产相关，也可能在同一资产跨期相关。双向聚类协方差可概念性写为
\[
\widehat V_{\rm two-way}=\widehat V_{\rm date}
+\widehat V_{\rm asset}-\widehat V_{\rm intersection}.
\]
聚类数太少时渐近近似不可靠，应报告簇数并考虑 wild bootstrap。标准误方法不能事后根据显著性选择。
